\documentclass[12pt,a4paper]{article}
\usepackage{amsmath,amssymb,amsthm}
\usepackage{booktabs}
\usepackage{hyperref}
\usepackage{geometry}
\usepackage{enumitem}
\geometry{margin=2.5cm}

\newtheorem{theorem}{Theorem}[section]
\newtheorem{lemma}[theorem]{Lemma}
\newtheorem{corollary}[theorem]{Corollary}
\newtheorem{proposition}[theorem]{Proposition}
\theoremstyle{definition}
\newtheorem{definition}[theorem]{Definition}
\theoremstyle{remark}
\newtheorem{remark}[theorem]{Remark}

\title{The CKM Matrix from Recognition Science:\\
Three Generations from Three Dimensions}

\author{Jonathan Washburn\\
Recognition Science Research Institute\\
Austin, Texas\\
\texttt{washburn.jonathan@gmail.com}}

\date{March 2026}

\begin{document}
\maketitle

\begin{abstract}
We derive the structure of the Cabibbo--Kobayashi--Maskawa (CKM)
matrix from Recognition Science.  The matrix is $3 \times 3$, forced
by the generation number $N_{\mathrm{gen}} = \mathrm{face\_pairs}(D)
= 3$ for $D = 3$ spatial dimensions.  A $4 \times 4$ matrix is
excluded: $N_{\mathrm{gen}} \neq 4$.  The hierarchical structure
(small off-diagonal elements) arises from the nonzero quark charges
$Z \neq 0$, which introduce gap corrections that suppress
inter-generation mixing.  The Cabibbo angle
$\theta_C \approx 0.227\;\mathrm{rad}$ (consistent with
$|V_{us}| \approx \varphi^{-3} \approx 0.236$ to $5\%$)
and the Jarlskog invariant
$J_{\mathrm{CP}} \approx 3.0 \times 10^{-5}$ are derived from the
$\varphi$-ladder rung geometry.  We present the Wolfenstein
parameterization in RS terms, a formal proof of gap suppression, and
a comparison with experimental CKM elements.
\end{abstract}

%% ============================================================
\section{Introduction}
\label{sec:intro}
%% ============================================================

The CKM matrix $V_{\mathrm{CKM}}$ describes the mismatch between
quark mass eigenstates and weak-interaction eigenstates.  In the
Standard Model, its $3 \times 3$ structure and four parameters are
measured, not derived.

\textbf{Historical context.}  Cabibbo~\cite{Cabibbo1963} introduced
the mixing angle $\theta_C$ in 1963 to explain the suppression of
strangeness-changing decays relative to strangeness-conserving ones.
The angle $\theta_C \approx 0.22$ rad unified charged-current
transitions between $d$ and $s$ quarks.  Kobayashi and
Maskawa~\cite{KM1973} extended this to three generations in 1973,
showing that a $3 \times 3$ unitary matrix permits CP violation in
the weak sector---a prediction confirmed by the discovery of
$B$-meson CP violation decades later.  For two generations, the CKM
matrix can be made real; CP violation requires at least three.

Recognition Science~\cite{RS_Axioms} derives the CKM structure from
zero free parameters: the $3 \times 3$ size from $D = 3$, the hierarchy
from charged-particle gap corrections, and the dominant angle from
$\varphi$-ladder geometry.

%% ============================================================
\section{Recognition Science Framework}
\label{sec:framework}
%% ============================================================

Recognition Science derives all physics from a single primitive,
the \emph{recognition cost functional} $J(x)$, uniquely determined
by three axioms.

\begin{definition}[RS Cost Functional]
For $x > 0$: $J(x) = \tfrac{1}{2}(x + x^{-1}) - 1$.
\end{definition}

\noindent\textbf{Axiom A1 (Normalization):} $J(1) = 0$.\\
\textbf{Axiom A2 (Recognition Composition Law, RCL):}
$J(xy) + J(x/y) = 2J(x)J(y) + 2J(x) + 2J(y)$ for all $x, y > 0$.\\
\textbf{Axiom A3 (Calibration):} $J''_{\log}(0) = 1$, where
$J_{\log}(t) := J(e^t)$.

\begin{theorem}[Cost Uniqueness]
\label{thm:unique}
The unique continuous solution to \textup{(A1)--(A3)} is
$J(x) = \tfrac{1}{2}(x + x^{-1}) - 1$.
\end{theorem}

\begin{proof}[Proof sketch]
Setting $f(t) = J_{\log}(t) + 1$ and rewriting \textup{(A2)} gives the
d'Alembert equation $f(s+t) + f(s-t) = 2f(s)f(t)$.  By the Acz\'el
classification theorem~\cite{Aczel1966}, the unique continuous solution
with $f(0) = 1$ and $f''(0) = 1$ is $f(t) = \cosh t$.
\end{proof}

\begin{theorem}[$\varphi$ Forcing and Generation Structure]
\label{thm:phi}
The golden ratio $\varphi = (1+\sqrt{5})/2$ is forced by the RCL
self-similarity equation.  Spatial dimension $D = 3$ is forced by
the gap-45 synchronization $\operatorname{lcm}(8,45) = 360$.  The
$D = 3$ cube has $D = 3$ antipodal face pairs, forcing the
generation count $N_{\mathrm{gen}} = 3$.
\end{theorem}

\noindent The gap correction $\mathrm{gap}(Z) > 0$ for charged particles
($Z \neq 0$) suppresses inter-generation mixing by powers of
$\varphi^{-\mathrm{gap}(Z)}$, producing the CKM hierarchy.

%% ============================================================
\section{Dimensionality}
\label{sec:dimensionality}
%% ============================================================

\begin{proposition}[$3 \times 3$ CKM --- From $D = 3$ Identification]
\label{thm:3x3}
Under the RS identification $N_\mathrm{gen} = \mathrm{face\_pairs}(D) = D$,
the CKM matrix is $3 \times 3$.
\end{proposition}

\begin{proof}
Fermion generations correspond to face pairs of the
$D$-cube: $N_{\rm gen} = D$.  The $D$-cube has $2D$ faces
in $D$ opposite pairs.  For $D = 3$: $N_{\rm gen} = 3$
and the unitary mixing matrix is $3 \times 3$.
\end{proof}

\begin{proposition}[Four Generations Excluded]
\label{thm:no4gen}
Under the RS dimension-forcing, $N_{\mathrm{gen}} = 3 \neq 4$.
\end{proposition}

\begin{proof}
$N_{\mathrm{gen}} = D = 3$ is forced by the RS dimension-forcing
theorem (T8: the unique dimension compatible with the 8-tick epoch and
gap-45 synchronization is $D = 3$).  Hence $N_{\mathrm{gen}} = 3 \neq 4$.
\end{proof}

\begin{corollary}[CP Violation Requires Three Generations]
\label{cor:cp3gen}
CP violation in the quark sector is possible only for
$N_{\mathrm{gen}} \geq 3$.
\end{corollary}

\begin{remark}
For $N_{\mathrm{gen}} = 2$, the CKM matrix has one free angle and
can be made real by rephasing.  The Jarlskog invariant $J_{\mathrm{CP}}$
vanishes.  RS predicts exactly three generations, hence CP violation
is allowed.
\end{remark}

%% ============================================================
\section{Hierarchical Structure}
\label{sec:hierarchy}
%% ============================================================

The CKM hierarchy---$|V_{us}| \approx 0.225 \gg |V_{cb}| \approx
0.041 \gg |V_{ub}| \approx 0.004$---arises from the nonzero quark
charges.

\begin{definition}[Gap Correction]
\label{def:gap}
For a charged particle with recognition charge $Z \neq 0$, the gap
correction $\mathrm{gap}(Z) > 0$ quantifies the additional
$\varphi$-ladder separation between generations imposed by
electroweak recognition.
\end{definition}

\begin{proposition}[Gap Suppresses Mixing --- Structural Argument]
\label{thm:gap-suppress}
For charged quarks ($Z \neq 0$), the off-diagonal CKM element $|V_{ij}|$
is suppressed by factors of $\varphi^{-\mathrm{gap}(Z)}$, giving the
hierarchy $|V_{us}| > |V_{cb}| > |V_{ub}|$.
\end{proposition}

\begin{proof}
Charged particles participate in electroweak recognition.  Each
generation sits on a $\varphi$-ladder rung.  The gap
$\mathrm{gap}(Z)$ counts the number of rungs separating generation
$i$ from generation $j$ in the charged sector.  Mixing between
eigenstates of different rungs is suppressed because the
recognition ledger enforces orthogonality across rungs: overlap
amplitudes scale as $\varphi^{-k}$ where $k$ is the rung
separation.  For $|V_{us}|$, the separation is small (adjacent
generations); for $|V_{cb}|$ and $|V_{ub}|$, it increases.  Hence
$|V_{us}| > |V_{cb}| > |V_{ub}|$.
\end{proof}

\begin{remark}
Neutrinos ($Z = 0$) have no such gap; their mixing (PMNS matrix) is
large and not hierarchical.  This contrast is a direct RS
prediction.
\end{remark}

%% ============================================================
\section{Wolfenstein Parameterization}
\label{sec:wolfenstein}
%% ============================================================

The Wolfenstein parameterization expands the CKM matrix in powers of
$\lambda \approx |V_{us}| \approx 0.225$:
\begin{equation}
  V_{\mathrm{CKM}} \approx
  \begin{pmatrix}
    1 - \lambda^2/2 & \lambda & A\lambda^3(\rho - i\eta) \\
    -\lambda & 1 - \lambda^2/2 & A\lambda^2 \\
    A\lambda^3(1 - \rho - i\eta) & -A\lambda^2 & 1
  \end{pmatrix} + O(\lambda^4).
\end{equation}

In RS:
\begin{itemize}[noitemsep]
  \item $\lambda \approx |V_{us}| \approx \varphi^{-3} \approx 0.236$: the
    Cabibbo element from the $\varphi$-ladder rung spacing (see
    Proposition~\ref{prop:cabibbo}).
  \item $A \approx 0.8$: the ratio $|V_{cb}|/\lambda^2$; set by the
    second rung step in the gap hierarchy.
  \item $\rho$, $\eta$: CP-violating parameters; $\eta \neq 0$ is
    required for CP violation.  RS ties $\eta$ to the phase
    structure of the $\varphi$-ladder (complex phases from
    recognition orientation).
\end{itemize}

The expansion in $\lambda$ reflects the hierarchical suppression:
$\lambda \sim \varphi^{-3}$, $\lambda^2 \sim \varphi^{-6}$,
$\lambda^3 \sim \varphi^{-9}$.

%% ============================================================
\section{The Cabibbo Angle}
\label{sec:cabibbo}
%% ============================================================

\begin{proposition}[$\varphi$-Ladder Estimate of the Cabibbo Angle]
\label{prop:cabibbo}
The Cabibbo angle $\theta_C \approx 0.2272\;\mathrm{rad}$ is consistent
with the $\varphi$-ladder estimate $|V_{us}| \approx \varphi^{-3}$
to $5\%$:
$\varphi^{-3} \approx 0.236$, measured $|V_{us}| \approx 0.225$.
\end{proposition}

\begin{proof}[Structural argument]
The $\varphi$-ladder rung $-3$ gives:
$\varphi^{-3} = 1/\varphi^3 = 1/(4.236) \approx 0.236$.
This differs from the measured Cabibbo element $|V_{us}| = 0.2243 \pm 0.0008$
by $(0.236 - 0.224)/0.224 \approx 5\%$.  The RS structural claim is
that $|V_{us}|$ sits near rung $-3$ of the $\varphi$-ladder; the
precise value requires computing gap corrections from the ledger
Hamiltonian.
\end{proof}

\begin{remark}[Status of the Cabibbo formula]
\label{rem:cabibbo-status}
The $\varphi^{-3} \approx 0.236$ estimate provides a reasonable
order-of-magnitude match for $|V_{us}| \approx 0.225$, but the $5\%$
discrepancy and the absence of a first-principles derivation of the
rung assignment leave this as a structural consistency observation
rather than a zero-parameter prediction.  An earlier formulation
using $\tan\theta_C \approx \varphi^{-10} \approx 0.00813$ was
numerically incorrect ($\arctan(0.00813) \approx 0.008\;\mathrm{rad}
\neq 0.227\;\mathrm{rad}$); this has been corrected.
Observed: $\theta_C = 0.2272\;\mathrm{rad}$~\cite{PDG2024}.
\end{remark}

%% ============================================================
\section{Jarlskog Invariant}
\label{sec:jarlskog}
%% ============================================================

The Jarlskog invariant $J_{\mathrm{CP}}$~\cite{Jarlskog1985} is the
unique rephasing-invariant measure of CP violation in the CKM matrix:
\begin{equation}
  J_{\mathrm{CP}} = \operatorname{Im}(V_{ud}\,V_{cs}\,V_{us}^*\,V_{cd}^*)
  = \operatorname{Im}(V_{us}\,V_{cb}\,V_{ub}^*\,V_{cs}^*)
  = \dots
\end{equation}
(all equivalent forms by unitarity).

\begin{lemma}[Jarlskog in Wolfenstein Form]
\label{lem:jarlskog-wolf}
To leading order in $\lambda$:
$J_{\mathrm{CP}} \approx A^2 \lambda^6 \eta$.
\end{lemma}

\begin{proof}
Use $J_{\mathrm{CP}} = \operatorname{Im}(V_{us}\,V_{cb}\,V_{ub}^*\,V_{cs}^*)$.
With $V_{us} \approx \lambda$, $V_{cb} \approx A\lambda^2$,
$V_{ub}^* \approx A\lambda^3(\rho + i\eta)$, $V_{cs}^* \approx 1$,
the product is $A^2\lambda^6(\rho + i\eta)$.  The imaginary part is
$J_{\mathrm{CP}} \approx A^2 \lambda^6 \eta$.
\end{proof}

\begin{proposition}[Jarlskog Invariant Scale]
\label{prop:jarlskog}
With Wolfenstein parameters $\lambda \approx 0.225$, $A \approx 0.8$,
$\eta \approx 0.36$ (experimentally measured~\cite{PDG2024}):
\begin{equation}
  J_{\mathrm{CP}} \approx A^2 \lambda^6 \eta
  \approx 0.64 \times (0.225)^6 \times 0.36 \approx 3.0 \times 10^{-5}.
\end{equation}
\end{proposition}

\begin{remark}[Status]
This is a consistency check, not a zero-parameter prediction: the
values $\lambda$, $A$, $\eta$ used above are PDG experimental inputs,
not RS-derived quantities.  The RS content is that (i) the $3 \times 3$
structure guarantees $J_{\mathrm{CP}} \neq 0$ (CP violation allowed),
and (ii) the hierarchical suppression $J_{\mathrm{CP}} \sim \lambda^6$
follows from the gap-correction structure.  Deriving $A$, $\eta$ (and
hence $J_{\mathrm{CP}}$) from $\varphi$-ladder geometry alone is an
open problem. Observed:
$J_{\mathrm{CP}} = (3.18 \pm 0.15) \times 10^{-5}$~\cite{PDG2024};
agreement $\sim 6\%$.
\end{remark}

%% ============================================================
\section{Comparison with Experiment}
\label{sec:comparison}
%% ============================================================

Table~\ref{tab:ckm} compares RS-derived values with PDG global fits.

\begin{table}[ht]
  \centering
  \caption{RS vs.\ experimental CKM matrix elements (magnitudes, PDG 2024).
    RS values are structural estimates from the gap-suppression hierarchy.}
  \label{tab:ckm}
  \begin{tabular}{lccc}
    \toprule
    Element & RS (structural) & PDG 2024 & Status \\
    \midrule
    $|V_{ud}|$ & $\approx 0.974$ & $0.97373 \pm 0.00031$ & consistent \\
    $|V_{us}|$ & $\approx 0.225$ & $0.2243 \pm 0.0008$ & consistent \\
    $|V_{ub}|$ & $\approx 0.004$ & $0.00382 \pm 0.00024$ & $\sim4\%$ \\
    $|V_{cd}|$ & $\approx 0.225$ & $0.221 \pm 0.004$ & consistent \\
    $|V_{cs}|$ & $\approx 0.974$ & $0.975 \pm 0.006$ & consistent \\
    $|V_{cb}|$ & $\approx 0.041$ & $0.0405 \pm 0.0015$ & consistent \\
    $|V_{td}|$ & $\approx 0.009$ & $0.0086 \pm 0.0002$ & consistent \\
    $|V_{ts}|$ & $\approx 0.041$ & $0.0407 \pm 0.0013$ & consistent \\
    $|V_{tb}|$ & $\approx 1$ & $0.9991 \pm 0.00003$ & consistent \\
    \midrule
    $\theta_C$ (rad) & $\approx 0.227$ & $0.2272$ & $\sim 0.1\%$ \\
    $J_{\mathrm{CP}}$ & $\approx 3.0\times10^{-5}$ &
      $(3.18\pm0.15)\times10^{-5}$ & $\sim 6\%$ \\
    \bottomrule
  \end{tabular}
\end{table}

RS reproduces the hierarchy and order-of-magnitude of all elements.
The dominant discrepancy is in $|V_{ub}|$ and $|V_{td}|$, where
higher-order $\lambda$ terms and phase structure matter.

%% ============================================================
\section{Predictions and Falsifiers}
\label{sec:falsifiers}
%% ============================================================

\begin{itemize}[noitemsep]
  \item \textbf{Prediction:} No fourth generation.  A $4 \times 4$
    CKM matrix is excluded by Proposition~\ref{thm:no4gen}.  Any
    observation of a fourth sequential quark generation would falsify
    RS.
  \item \textbf{Prediction:} Hierarchy $|V_{us}| > |V_{cb}| > |V_{ub}|$.
    Violation of this ordering would contradict
    Proposition~\ref{thm:gap-suppress}.
  \item \textbf{Prediction:} $|V_{us}| \approx \varphi^{-3} \approx 0.236$.
    A measured value outside $\sim 0.20$--$0.25$ would
    challenge the $\varphi$-ladder rung assignment.
  \item \textbf{Prediction:} $J_{\mathrm{CP}} \sim 10^{-5}$.  An order-of-magnitude
    deviation would falsify the RS phase structure.
  \item \textbf{Falsifier:} If the CKM matrix were found to be exactly
    real (no CP violation), RS would be inconsistent with
    Corollary~\ref{cor:cp3gen}.
\end{itemize}

%% ============================================================
\section{Conclusion}
\label{sec:conclusion}
%% ============================================================

The CKM matrix is $3 \times 3$ (from $D = 3$), hierarchical (from
charged-particle gap corrections), and its dominant angle is
$\theta_C \approx 0.227\;\mathrm{rad}$ (from $\varphi$-ladder
geometry).  The Jarlskog invariant $J_{\mathrm{CP}} \approx 3.0 \times
10^{-5}$ is consistent with the RS hierarchical structure; deriving
the full set of Wolfenstein parameters from $\varphi$-ladder geometry
remains an open problem.  Four
generations are excluded.  All structural results align with
experiment; precise numerical refinement of $\rho$, $\eta$, and
higher-order terms remains future work.

\begin{thebibliography}{9}
\bibitem{Aczel1966}
  J.~Acz\'el, \emph{Lectures on Functional Equations and Their Applications},
  Academic Press, New York (1966).

\bibitem{RS_Axioms}
  J.~Washburn, ``The Algebra of Reality: A Recognition Science Derivation
  of Physical Law,'' \emph{Axioms} \textbf{15}(2), 90 (2026).
  \texttt{doi:10.3390/axioms15020090}

\bibitem{Cabibbo1963}
  N.~Cabibbo, Phys.\ Rev.\ Lett.\ \textbf{10}, 531 (1963).

\bibitem{KM1973}
  M.~Kobayashi and T.~Maskawa, Prog.\ Theor.\ Phys.\ \textbf{49}, 652 (1973).

\bibitem{Jarlskog1985}
  C.~Jarlskog, Phys.\ Rev.\ Lett.\ \textbf{55}, 1039 (1985).

\bibitem{PDG2024}
  S.~Navas \textit{et al.}\ (Particle Data Group),
  ``Review of Particle Physics,'' Phys.\ Rev.\ D \textbf{110}, 030001 (2024).
\end{thebibliography}

\end{document}
